\section*{Bab \ref{ch:g11:becoming-male-female} --- Menjadi Lelaki atau Perempuan}

\begin{solution}{exo:g11:becoming-male-female:1}
Gamet ayahnya: tiap \omterm{def:g10:cells-common-unit:cell}{sel} telur membawa satu X, sedangkan separuh
spermanya membawa X dan separuhnya Y, sehingga XX dan XY sama mungkinnya.
\end{solution}

\begin{solution}{exo:g11:becoming-male-female:2}
Dua gonad yang belum terbedakan, yang masing-masing punya dua saluran
di sampingnya: saluran Wolff (bakal saluran jantan) dan saluran Müller
(bakal saluran betina). Serupa pada embrio XX dan XY.
\end{solution}

\begin{solution}{exo:g11:becoming-male-female:3}
Pada \omterm{def:g11:cell-cycle-mitosis:chromosome}{kromosom} Y; terungkapkan dalam gonad yang belum terbedakan sekitar
pekan ke-7, dan \omterm{def:g11:gene-expression:protein}{proteinnya} menyalakan \omterm{def:g10:universal-dna:gene}{gen} yang membuat gonadnya menjadi
sebuah testis.
\end{solution}

\begin{solution}{exo:g11:becoming-male-female:4}
\omterm{prop:g11:becoming-male-female:hormones}{Testosteron} menjaga dan mengembangkan saluran Wolffnya menjadi saluran
lelaki; hormon anti-Müller membuat saluran Müllernya mengerut.
\end{solution}

\begin{solution}{exo:g11:becoming-male-female:5}
Sebuah pusat di otaknya melepaskan denyut sebuah hormon yang membuat
hipofisisnya mengeluarkan LH dan FSH, yang lalu mengaktifkan gonadnya.
Ciri kelamin sekunder: perubahan yang dihasilkan hormon kelaminnya ---
rambut tubuh, suara, payudara, pembagian otot dan lemaknya, pacu
tumbuhnya, pertumbuhan organ perkembangbiakannya.
\end{solution}

\begin{solution}{exo:g11:becoming-male-female:6}
Janin itu mengembangkan saluran betina, meskipun punya \omterm{def:g11:cell-cycle-mitosis:chromosome}{kromosom} Y:
saluran jantannya menuntut hormon testisnya, dan saluran betinanya
terbentuk tanpa gonad apa pun.
\end{solution}

\begin{solution}{exo:g11:becoming-male-female:7}
\omterm{prop:g11:becoming-male-female:hormones}{Testosteron} membangun saluran jantannya tetapi tidak menyingkirkan
saluran betinanya; testisnya pasti mengeluarkan zat kedua untuk itu ---
yaitu AMH.
\end{solution}

\begin{solution}{exo:g11:becoming-male-female:8}
Testis (Sry bekerja); saluran Wolffnya terjaga, saluran Müllernya
mengerut; anatomi luarnya jantan. Mandul: penghasilan spermanya
menuntut \omterm{def:g10:universal-dna:gene}{gen} Y yang tidak dipasok transgennya.
\end{solution}

\begin{solution}{exo:g11:becoming-male-female:9}
Estrogen: 50\% pada sekitar umur 11, dataran pada sekitar 14--15.
\omterm{prop:g11:becoming-male-female:hormones}{Testosteron}: 50\% pada sekitar umur 13, dataran pada sekitar 16--17.
\end{solution}

\begin{solution}{exo:g11:becoming-male-female:10}
Tulangnya berhenti tumbuh ketika hormon kelaminnya mencapai aras
dewasanya; estrogen mencapainya sekitar dua tahun sebelum testosteron,
sehingga lempeng pertumbuhan anak perempuan menutup lebih awal.
\end{solution}

\begin{solution}{exo:g11:becoming-male-female:11}
Tanpa testis: gonadnya menjadi gonad yang menyerupai ovarium; tanpa
testosteron atau AMH: saluran Müllernya berkembang, saluran Wolffnya
mengerut, dan anatomi luarnya perempuan saat lahir. Pada \omterm{prop:g11:becoming-male-female:puberty}{pubertasnya}
gonadnya, yang tidak punya \omterm{def:g10:cells-common-unit:cell}{sel} telur yang normal, mengeluarkan sedikit
saja; tanpa pengobatan hormon, ciri sekundernya berkembang buruk.
\end{solution}

\begin{solution}{exo:g11:becoming-male-female:12}
Testis (\omterm{prop:g11:becoming-male-female:sry}{SRY} bekerja). Saluran Wolffnya mengerut (tanpa tanggapan
terhadap testosteron); saluran Müllernya mengerut (AMH bekerja):
sehingga tidak ada saluran dalam dari jenis mana pun. Anatomi luarnya
perempuan (bentuk lelakinya menuntut kerja testosteron). Pada
\omterm{prop:g11:becoming-male-female:puberty}{pubertasnya}, testosteron yang diubah menjadi estrogen menghasilkan
payudara dan perawakan perempuan, tanpa rambut tubuh.
\end{solution}

\begin{solution}{exo:g11:becoming-male-female:13}
Begitu sebuah gonad menjadi testis atau ovarium, hormonnya dapat
membentuk saluran dan tubuhnya persis seperti pada mamalia; yang
berbeda adalah langkah pertamanya --- sebuah proses yang peka suhu
menggantikan \omterm{prop:g11:becoming-male-female:sry}{SRY} sebagai sakelar yang memutuskan gonadnya.
\end{solution}

\begin{solution}{exo:g11:becoming-male-female:14}
\omterm{def:g11:cell-cycle-mitosis:chromosome}{Kromosom} X membawa ratusan \omterm{def:g10:universal-dna:gene}{gen} yang diperlukan setiap \omterm{def:g10:cells-common-unit:cell}{sel}, sehingga
tidak ada embrio yang bertahan tanpa satu X; \omterm{def:g11:cell-cycle-mitosis:chromosome}{kromosom} Y hampir tidak
membawa apa pun selain \omterm{def:g10:universal-dna:gene}{gen} penentu testis dan \omterm{def:g10:universal-dna:gene}{gen} yang berkaitan dengan
sperma, sehingga individu XX tidak kehilangan apa pun yang
diperlukannya. Sebuah Y hanya berpindah kepada anak lelaki karena anak
perempuannya menerima \omterm{def:g11:cell-cycle-mitosis:chromosome}{kromosom} X ayahnya.
\end{solution}

\begin{solution}{exo:g11:becoming-male-female:15}
Ia berarti bahwa saluran perempuannya berkembang ketika tidak ada
hormon testis yang turun tangan --- yaitu kenyataan tentang isyarat
mana yang diperlukan, bukan sebuah peringkat. Ia mencerminkan rencana
embrio yang bersama yang di atasnya satu garis keturunan \omterm{def:g10:cells-common-unit:cell}{sel}, yaitu
testisnya, memaksakan pilihan lain lewat dua hormon: kedua kelaminnya
adalah dua hasil satu program, yang berbeda oleh sebuah sakelar.
\end{solution}

\begin{solution}{pb:g11:becoming-male-female:1}
\textbf{1.} Ovariumnya tidak diperlukan bagi saluran betinanya, yang
terbentuk tanpa gonad apa pun; testisnya diperlukan bagi saluran
jantannya dan bagi penyingkiran saluran betinanya.

\textbf{2.} Setelah hari ke-22 salurannya sudah mulai terbedakan di
bawah pengaruh gonadnya; percobaannya harus bekerja sebelum
keputusannya diambil.

\textbf{3.} Hormon testisnya, yang dibawa darah ke salurannya:
cangkoknya, yang jauh dari salurannya, tetap menjantankannya.

\textbf{4.} \omterm{prop:g11:becoming-male-female:hormones}{Testosteron} membangun saluran jantannya; ia tidak
menyingkirkan saluran betinanya.

\textbf{5.} Zat testis yang kedua pasti membuat saluran Müllernya
mengerut; keberadaannya dapat dibenarkan dengan menyarinya dari testis
dan menanamkannya sendirian, dengan mengharapkan saluran betinanya
lenyap dan saluran jantannya tidak terbentuk --- dan itulah yang
dikerjakan AMH.

\textbf{6.} Pada sisi yang dioperasi, tanpa testosteron atau AMH:
saluran betina, tanpa saluran jantan. Pada sisi yang utuh: saluran
jantan, saluran betinanya mengerut.

\textbf{7.} Pada sisi cangkoknya, saluran jantan dan tanpa saluran
betina; pada sisi yang lain, hanya saluran betina.

\textbf{8.} Pada sisi kristalnya, kedua salurannya hadir; pada sisi
yang lain, hanya saluran betinanya.

\textbf{9.} Tiap janinnya bertugas sebagai kontrolnya sendiri: beda
antara kedua sisinya hanya dapat disebabkan kehadiran setempat
testisnya, testosteronnya atau ketiadaannya.

\textbf{10.} Ovariumnya disingkirkan atau dikalahkan dan testis yang
dicangkokkan, yang berasal dari hewan lain, tidak membuat sperma dalam
inang yang asing; salurannya jantan tetapi tidak ada gamet yang dihasilkan.

\textbf{11.} \omterm{prop:g11:becoming-male-female:sry}{SRY} hadir: testis; testosteron dan AMH: saluran lelaki,
anatomi lelaki; pada \omterm{prop:g11:becoming-male-female:puberty}{pubertasnya}, testosteron: ciri lelaki. Mandul
(tanpa \omterm{def:g10:universal-dna:gene}{gen} Y bagi spermanya).

\textbf{12.} Tanpa \omterm{prop:g11:becoming-male-female:sry}{SRY} yang bekerja: tanpa testis; tanpa hormonnya:
saluran dan anatomi perempuan; gonadnya berkembang buruk.

\textbf{13.} Testis; saluran Wolffnya berkembang (testosteron bekerja);
saluran Müllernya ikut bertahan (tanpa tanggapan terhadap AMH):
seorang lelaki dengan, sebagai tambahan, sebuah rahim dan saluran
telur; anatomi dan \omterm{prop:g11:becoming-male-female:puberty}{pubertasnya} lelaki.

\textbf{14.} Hanya yang ketiga: testis dan saluran lelakinya berfungsi,
walaupun rahim yang bertahan itu dapat menghalangi turunnya testisnya.
Yang pertama tidak punya \omterm{def:g10:universal-dna:gene}{gen} spermanya; yang kedua tidak punya gonad
yang berfungsi.

\textbf{15.} Kelamin \omterm{def:g11:cell-cycle-mitosis:chromosome}{kromosomnya} meramalkan kelamin anatominya hanya
lewat langkah di antaranya: \omterm{prop:g11:becoming-male-female:sry}{SRY} yang bekerja, hormon yang bekerja,
reseptor yang bekerja. Ubahlah langkah mana pun dan kedua arasnya berpisah.

\textbf{16.} Anak perempuan: 20\% pada umur 10, pacu tumbuhnya sejak
sekitar 11, berhenti pada sekitar 15. Anak lelaki: 20\% pada sekitar
12.5, pacu tumbuhnya sejak sekitar 13.5, berhenti pada sekitar 17.

\textbf{17.} Anak perempuan sekitar 4 tahun, anak lelaki sekitar 3.5
--- tetapi anak lelakinya berangkat dari alas yang lebih tinggi, karena
telah tumbuh pada laju masa kanak-kanaknya dua tahun lebih lama.

\textbf{18.} Anak lelaki sekitar $3.5 \times 8 = \qty{28}{cm}$, anak
perempuan $4 \times
7 = \qty{28}{cm}$: pacu tumbuhnya serupa; beda dewasanya sebagian besar
berasal dari dua tahun tambahan pertumbuhan masa kanak-kanak sebelum
pacu tumbuh anak lelakinya.

\textbf{19.} Kenaikan yang dini membawa lempeng pertumbuhannya menutup
bertahun-tahun lebih awal: anaknya tumbuh cepat sekarang tetapi
berhenti pada tinggi dewasa yang pendek. Menunda hormonnya
memperpanjang pertumbuhan masa kanak-kanaknya.

\textbf{20.} \omterm{prop:g11:becoming-male-female:hormones}{Testosteron}, yang membangun saluran lelakinya, dan hormon
anti-Müller, yang menyingkirkan saluran perempuannya; keduanya dibuat
testisnya, yaitu organ yang dinyalakan satu \omterm{def:g10:universal-dna:gene}{gen} tunggal \omterm{prop:g11:becoming-male-female:sry}{SRY}.
\end{solution}
